% ============================================================
% FASE 1 DEL HITO 5: INSTRUMENTACIÓN Y TRAZABILIDAD
% Bloque listo para integrar en la memoria del proyecto VinOps
% ============================================================

\section{Hito 5: Monitorización y Feedback Loop}

\subsection{Fase 1 -- Instrumentación y Trazabilidad}

\subsubsection{Objetivo y alcance}

El objetivo de esta fase es dotar al sistema VinOps de capacidades de observabilidad que permitan registrar, en tiempo real, tanto métricas operativas del servicio como métricas proxy de calidad del modelo. Esta instrumentación constituye la base sobre la que se construirán las fases posteriores de visualización (Prometheus/Grafana), detección de drift y feedback loop.

\subsubsection{Diseño de la arquitectura de observabilidad}

Se ha diseñado una arquitectura de tres capas:

\begin{enumerate}
    \item \textbf{Capa de instrumentación:} integrada directamente en el servicio de inferencia, emite métricas en cada petición.
    \item \textbf{Capa de persistencia:} logs estructurados en formato JSONL y métricas Prometheus expuestas vía HTTP.
    \item \textbf{Capa de consumo:} endpoints de consulta (\texttt{/health}, \texttt{/monitor}, \texttt{/metrics}) y ficheros de log accesibles desde el host.
\end{enumerate}

\subsubsection{Métricas operativas}

Siguiendo los requisitos del profesorado, se han instrumentado las siguientes métricas operativas mediante la librería \texttt{prometheus\_client}:

\begin{table}[htbp]
\centering
\caption{Métricas operativas instrumentadas en VinOps}
\label{tab:metricas-operativas}
\begin{tabular}{@{}llp{6cm}@{}}
\toprule
\textbf{Métrica} & \textbf{Tipo} & \textbf{Descripción} \\
\midrule
\texttt{vinops\_requests\_total} & Counter & Peticiones HTTP por método, endpoint y código de estado \\
\texttt{vinops\_request\_latency\_seconds} & Histogram & Latencia de las peticiones a \texttt{/predict} (buckets configurados) \\
\texttt{vinops\_cpu\_usage\_percent} & Gauge & Uso de CPU del proceso de inferencia \\
\texttt{vinops\_memory\_usage\_bytes} & Gauge & Consumo de RAM residente del proceso \\
\bottomrule
\end{tabular}
\end{table}

Estas métricas se actualizan en cada petición y se exponen en el endpoint \texttt{GET /metrics} en formato texto compatible con Prometheus.

\subsubsection{Métricas proxy del modelo}

Dado que en producción la etiqueta real del cultivar no está disponible en el momento de la predicción, se ha optado por un conjunto de métricas \textit{proxy} que detectan degradación del modelo sin requerir \textit{ground truth}:

\begin{table}[htbp]
\centering
\caption{Métricas proxy del modelo implementadas}
\label{tab:metricas-proxy}
\begin{tabular}{@{}llp{6.5cm}@{}}
\toprule
\textbf{Métrica} & \textbf{Tipo} & \textbf{Justificación enológica} \\
\midrule
\texttt{vinops\_prediction\_confidence} & Histogram & Confianza máxima de la predicción. Si baja respecto al período de entrenamiento, el modelo está en zona gris (frontera entre cultivares). \\
\texttt{vinops\_prediction\_entropy} & Histogram & Entropía de la distribución de probabilidades. Entropía alta indica confusión entre clases, típica del solapamiento clase~2--clase~3 detectado en el Hito~3. \\
\texttt{vinops\_predictions\_by\_class\_total} & Counter & Distribución de clases predichas. Desviaciones respecto a la distribución del entrenamiento (33\%/40\%/27\%) alertan de cambios en la línea de producción. \\
\texttt{vinops\_anomaly\_flags\_total} & Counter & Muestras cuya composición fisicoquímica supera $\pm 3\sigma$ respecto al dataset de referencia (umbrales del Hito~2). \\
\texttt{vinops\_low\_confidence\_total} & Counter & Predicciones con confianza $<0.70$, umbral que marca el límite de fiabilidad operativa. \\
\bottomrule
\end{tabular}
\end{table}

\subsubsection{Detección de anomalías en inferencia}

Se ha reutilizado la metodología del Hito~2 adaptándola a tiempo real. Para cada muestra entrante se calcula el z-score de cada variable respecto a la media y desviación del dataset de entrenamiento:

\begin{equation}
    z_j = \frac{x_j - \bar{x}_j}{\sigma_j}, \qquad j = 1,\dots,13
\end{equation}

Si $|z_j| > 3$ para cualquier variable, la muestra se marca como anómala. Esta información se incluye tanto en la respuesta JSON de \texttt{/predict} como en el log estructurado, permitiendo al enólogo revisar las variables específicas que desvían del perfil esperado.

\subsubsection{Logging estructurado}

Cada predicción genera una línea JSON en el fichero \texttt{logs/vinops\_predictions.jsonl} con la siguiente estructura:

\begin{itemize}
    \item \texttt{timestamp}: instante ISO 8601 de la predicción.
    \item \texttt{input}: vector completo de las 13 variables fisicoquímicas.
    \item \texttt{prediction} y \texttt{probabilities}: clase predicha y distribución de probabilidades.
    \item \texttt{confidence} y \texttt{entropy}: métricas proxy de calidad.
    \item \texttt{anomaly}: resultado del detector con z-scores por variable.
    \item \texttt{latency\_ms}: tiempo de respuesta de la inferencia.
\end{itemize}

Este formato garantiza trazabilidad completa y facilita análisis posteriores mediante herramientas de procesamiento de logs.

\subsubsection{Endpoints del servicio monitorizado}

\begin{table}[htbp]
\centering
\caption{Endpoints de la API instrumentada}
\label{tab:endpoints}
\begin{tabular}{@{}llp{7cm}@{}}
\toprule
\textbf{Endpoint} & \textbf{Método} & \textbf{Descripción} \\
\midrule
\texttt{/health} & GET & Estado del servicio, metadatos del modelo, resumen de monitorización y uso de CPU/RAM \\
\texttt{/predict} & POST & Predicción + métricas proxy + detección de anomalías + logging \\
\texttt{/metrics} & GET & Métricas Prometheus en formato text/plain \\
\texttt{/monitor} & GET & Resumen agregado: predicciones totales, anomalías, confianza media, distribución de clases \\
\bottomrule
\end{tabular}
\end{table}

\subsubsection{Evidencias de funcionamiento}

A continuación se presentan los resultados de las pruebas ejecutadas sobre el servicio desplegado.

\paragraph{Predicción normal.} Muestra típica de la clase~1. El modelo predice con confianza 0.995 y entropía 0.045. Ninguna variable supera el umbral de anomalía.

\paragraph{Predicción de frontera.} Muestra con valores intermedios. El modelo predice clase~2 con confianza 0.79 y entropía 0.91, indicando mayor incertidumbre pero dentro de rangos normales.

\paragraph{Predicción anómala.} Muestra con valores extremos pero dentro de los rangos físicos del vino. El detector marca anomalía en \texttt{Ash}, \texttt{Magnesium} y \texttt{Proanthocyanins} (z-scores $>3$). La confianza del modelo baja a 0.755 y la entropía sube a 0.99, señalando que la muestra cae fuera del perfil aprendido.

\paragraph{Endpoint \texttt{/metrics}.} Se verifica la exposición correcta de métricas Prometheus, incluyendo contadores de peticiones, histogramas de latencia y gauges de CPU/RAM.

\paragraph{Endpoint \texttt{/monitor}.} Tras las tres predicciones de prueba, el resumen agregado reporta: 3 predicciones totales, 1 anomalía detectada, confianza media 0.847 y distribución equilibrada entre las tres clases.

\paragraph{Logs estructurados.} Se verifica la generación de 3 líneas JSON en \texttt{logs/vinops\_predictions.jsonl}, cada una con la trazabilidad completa descrita.

\subsubsection{Decisiones de diseño tomadas}

\begin{enumerate}
    \item \textbf{Separación en carpeta \texttt{observabilidad/}:} se mantiene desacoplada del código del Hito~4 para no contaminar el despliegue base, facilitando la revisión independiente por el tribunal.
    \item \textbf{Uso de métricas proxy en lugar de F1/AUC online:} se justifica porque la etiqueta real no está disponible en el momento de la predicción. El F1-Score se recalculará offline en la Fase~3 (feedback loop).
    \item \textbf{Reutilización de umbrales $\pm 3\sigma$ del Hito~2:} garantiza coherencia metodológica entre el análisis exploratorio y la monitorización operativa.
    \item \textbf{Validación de entrada con Pydantic:} rechaza valores físicamente imposibles (p.\ ej.\ alcohol $>16\%$) antes de la inferencia, actuando como primera línea de defensa.
\end{enumerate}

\subsubsection{Próximas fases}

\begin{itemize}
    \item \textbf{Fase 2:} Despliegue de Prometheus + Grafana para visualización de métricas.
    \item \textbf{Fase 3:} Endpoint \texttt{/feedback} y retraining semi-automático.
    \item \textbf{Fase 4:} Integración con MLflow Tracking del bonus del Hito~4.
\end{itemize}
